% TikZ 图：传感器分压模型
% Day 19：恒定电流与基本电路
\begin{tikzpicture}[scale=0.92,>=Latex,line cap=round,line join=round]
  \path[use as bounding box] (-4.55,-2.15) rectangle (4.55,2.10);

  \tikzset{
    wire/.style={thick},
    res/.style={thick,fill=white},
    signal/.style={->,very thick,blue},
    current/.style={->,very thick,red}
  }

  % Circuit frame: keep the circuit uncluttered; explanations are placed below.
  \draw[wire] (-3.10,1.30) -- (-1.55,1.30);
  \draw[res] (-1.55,1.00) rectangle (-0.35,1.60);
  \node at (-0.95,1.30) {$R_0$};
  \draw[wire] (-0.35,1.30) -- (1.45,1.30);
  \draw[res] (1.45,0.18) rectangle (2.05,1.30);
  \node at (1.75,0.74) {$R_x$};
  \draw[wire] (1.75,0.18) -- (1.75,-0.55) -- (-3.10,-0.55);
  \draw[wire] (-3.10,-0.55) -- (-3.10,-0.38);
  \draw[wire] (-3.10,0.28) -- (-3.10,1.30);

  % Battery
  \draw[thick] (-3.45,0.28) -- (-2.75,0.28);
  \draw[thick] (-3.30,-0.38) -- (-2.90,-0.38);
  \node[left] at (-3.48,0.28) {$+$};
  \node[left] at (-3.48,-0.38) {$-$};
  \node[left] at (-3.52,-0.05) {$E$};

  % Current and output voltage
  \draw[current] (-2.25,1.52) -- (-1.72,1.52);
  \node[red,font=\small] at (-2.00,1.82) {$I$};
  \draw[signal] (2.65,0.18) -- (2.65,1.30);
  \draw[wire] (1.75,1.30) -- (2.75,1.30);
  \draw[wire] (1.75,0.18) -- (2.75,0.18);
  \node[blue,right] at (2.70,0.74) {$U_{\text{出}}$};

  % Input physical quantities
  \node[font=\small,align=center] at (0,-1.20) {温度、光照、压力等改变 $R_x$};
  \node[font=\small,align=center] at (0,-1.65) {物理量变化 $\to$ 电阻变化 $\to$ 输出电压变化};
\end{tikzpicture}
